\section{Introduction}
\label{sec:intro}

\begin{figure*}[t]
\centering
\resizebox{\textwidth}{!}{% Overview figure (TikZ). Included by main.tex; needs tikz with the libraries arrows.meta, calc, patterns and
% decorations.pathreplacing. Colours: Okabe-Ito (colour-blind safe), as in the data figures.
\definecolor{oiBlue}{HTML}{0072B2}\definecolor{oiOrange}{HTML}{E69F00}\definecolor{oiVerm}{HTML}{D55E00}
\definecolor{oiGreen}{HTML}{009E73}\definecolor{oiPink}{HTML}{CC79A7}
\begin{tikzpicture}[
    x=1cm, y=1cm, font=\scriptsize, >={Stealth[length=1.6mm]},
    panel/.style={draw=#1, line width=0.6pt, rounded corners=2pt},
    head/.style={fill=#1, text=white, font=\footnotesize\bfseries, anchor=north west, inner sep=2pt,
                 minimum height=3.6mm, rounded corners=1pt},
    robot/.style={circle, fill=#1, draw=black!70, line width=0.3pt, minimum size=2.4mm, inner sep=0pt},
    task/.style={circle, draw=black!80, fill=white, line width=0.5pt, minimum size=4.2mm, inner sep=0pt,
                 font=\tiny\bfseries},
    depot/.style={rectangle, fill=#1, draw=black!70, line width=0.3pt, minimum size=2.4mm, inner sep=0pt},
    tbox/.style={rectangle, rounded corners=1pt, draw=black!70, fill=black!4, minimum width=6.2mm,
                 minimum height=4.2mm, inner sep=1pt, font=\tiny\bfseries},
    step/.style={rectangle, rounded corners=4pt, draw=#1, fill=#1!12, line width=0.5pt, align=center,
                 inner sep=1.2pt, font=\tiny, text width=10.5mm, minimum height=6.2mm},
    meth/.style={rectangle, rounded corners=1pt, draw=black!50, fill=black!3, inner sep=1pt,
                 minimum width=21mm, minimum height=3.1mm, font=\tiny, anchor=west},
    chev/.style={rectangle, rounded corners=1pt, draw=#1, fill=#1!12, font=\tiny, align=center, inner sep=1.2pt,
                 minimum height=6.5mm, text width=9.2mm},
    note/.style={font=\tiny, align=left, anchor=west, execute at begin node={\hyphenpenalty=10000}},
  ]
% ---------------------------------------------------------------------------------------------------------------
% (a) Problem
\draw[panel=oiBlue] (0,0) rectangle (5.25,5.3);
\node[head=oiBlue] at (0.05,5.25) {(a) Coalitions with synchronised starts};
\node[depot=oiBlue] (dA) at (0.4,4.35) {};
\node[robot=oiBlue] (r1) at (0.95,4.35) {};
\node[anchor=south west, font=\tiny, text=oiBlue, inner sep=1pt] at ([xshift=-1.5mm]r1.north) {traits $a_1{=}(1,0,1)$};
\node[depot=oiOrange] (dB) at (0.4,2.7) {};
\node[robot=oiOrange] (r2) at (0.95,2.7) {};
\node[anchor=north west, font=\tiny, text=oiOrange!80!black, inner sep=1pt] at ([xshift=-1.5mm]r2.south)
  {traits $a_2{=}(0,1,1)$};
\node[task] (t1) at (2.75,3.5) {1};
\node[note, align=left] at ([xshift=0.4mm]t1.east) {needs\\$r_1{=}(1,1,0)$};
\node[task] (t2) at (4.7,4.5) {2};
\node[task] (t3) at (4.7,2.55) {3};
\draw[->, dashed, oiBlue, line width=0.5pt] (r1) to[bend left=10] (t1);
\draw[->, dashed, oiOrange!85!black, line width=0.5pt] (r2) to[bend right=10] (t1);
\draw[->, dashed, oiBlue, line width=0.5pt] (t1) to[bend left=12] (t2);
\draw[->, dashed, oiOrange!85!black, line width=0.5pt] (t1) to[bend right=12] (t3);
\node[depot=black!45] at (0.35,2.05) {};
\node[note] at (0.45,2.05) {depot};
\node[robot=black!45] at (1.25,2.05) {};
\node[note] at (1.35,2.05) {robot};
\node[task, minimum size=2.4mm] at (2.1,2.05) {};
\node[note] at (2.2,2.05) {task};
\node[note] at (2.9,2.05) {goal: min.\ makespan};
% timeline of task 1
\draw[->, black!70] (0.35,0.45) -- (5.1,0.45) node[anchor=north east, font=\tiny] {time};
\node[anchor=east, font=\tiny, text=oiBlue] at (0.98,1.38) {robot 1};
\node[anchor=east, font=\tiny, text=oiOrange!80!black] at (0.98,0.88) {robot 2};
\draw[oiBlue, line width=0.5pt] (1.0,1.38) -- (1.7,1.38);
\fill[oiBlue] (1.7,1.38) circle (0.6mm);
\fill[pattern=north east lines, pattern color=oiBlue!70] (1.7,1.28) rectangle (2.75,1.48);
\node[anchor=south, font=\tiny] at (2.22,1.48) {arrives, waits};
\draw[oiOrange, line width=0.5pt] (1.0,0.88) -- (2.75,0.88);
\fill[oiOrange] (2.75,0.88) circle (0.6mm);
\fill[oiGreen!55] (2.75,0.78) rectangle (3.85,1.48);
\node[font=\tiny, align=center] at (3.3,1.13) {task 1\\together};
\draw[densely dotted] (2.75,0.45) -- (2.75,1.62);
\draw[densely dotted] (3.85,0.45) -- (3.85,1.62);
\node[anchor=north, font=\tiny] at (2.75,0.45) {$s_1$};
\node[anchor=north, font=\tiny] at (3.85,0.45) {$f_1{=}s_1{+}d_1$};
\node[note] at (3.92,1.13) {on to 2, 3};
% ---------------------------------------------------------------------------------------------------------------
% (b) Method
\draw[panel=oiVerm] (5.45,0) rectangle (11.85,5.3);
\node[head=oiVerm] at (5.5,5.25) {(b) ALNS over a deadlock-free plan};
\node[note] at (5.55,4.68) {plan: a global task order $\pi$ and a minimal cover $C_j$ per task};
\node[tbox] (p1) at (6.0,4.12) {1};
\node[tbox] (p2) at (7.0,4.12) {3};
\node[tbox] (p3) at (8.0,4.12) {2};
\node[tbox] (p4) at (9.0,4.12) {$\cdots$};
\draw[->] (p1) -- (p2); \draw[->] (p2) -- (p3); \draw[->] (p3) -- (p4);
\node[robot=oiBlue, minimum size=1.8mm] at ([yshift=-3.3mm, xshift=-1.2mm]p1.south) {};
\node[robot=oiOrange, minimum size=1.8mm] at ([yshift=-3.3mm, xshift=1.2mm]p1.south) {};
\node[robot=oiOrange, minimum size=1.8mm] at ([yshift=-3.3mm]p2.south) {};
\node[robot=oiBlue, minimum size=1.8mm] at ([yshift=-3.3mm]p3.south) {};
\node[note, text width=22mm] at (9.45,4.0) {robots visit their tasks in the order $\pi$: no deadlock; forward pass = simulator};
% ALNS loop
\coordinate (c) at (7.05,1.95);
\node[step=oiBlue] (d) at ($(c)+(0,0.95)$) {destroy\\$q$ tasks};
\node[step=oiOrange] (r) at ($(c)+(1.0,0)$) {exact\\repair};
\node[step=oiPink] (a) at ($(c)+(0,-0.95)$) {accept by\\annealing};
\node[step=oiGreen] (w) at ($(c)+(-1.0,0)$) {adapt\\weights};
\draw[->, line width=0.6pt] (d.east) to[bend left=35] (r.north);
\draw[->, line width=0.6pt] (r.south) to[bend left=35] (a.east);
\draw[->, line width=0.6pt] (a.west) to[bend left=35] (w.south);
\draw[->, line width=0.6pt] (w.north) to[bend left=35] (d.west);
\node[font=\footnotesize\bfseries] at (c) {ALNS};
\node[font=\tiny] at (7.05,0.22) {$\times 8$ workers share their best plan every 1\,s};
% exact insertion inset
\draw[rounded corners=2pt, draw=black!60, fill=black!3] (8.7,0.15) rectangle (11.75,3.35);
\node[anchor=north west, font=\tiny\bfseries] at (8.75,3.32) {exact insertion of task $j$ in $O(|C_j|)$};
\node[tbox, minimum width=4.5mm] (q1) at (9.15,2.45) {$p$};
\node[tbox, minimum width=4.5mm, draw=oiOrange, fill=oiOrange!15] (qj) at (10.0,2.45) {$j$};
\node[tbox, minimum width=4.5mm] (qk) at (10.85,2.45) {$k_i$};
\draw[->] (q1) -- (qj); \draw[->] (qj) -- (qk);
\draw[->, black!45, densely dashed] (q1.north) to[bend left=28] (qk.north);
\draw[decorate, decoration={brace, amplitude=2pt, mirror}] ([yshift=-1mm]qk.south west) --
  ([yshift=-1mm, xshift=5mm]qk.south east) node[midway, below=1pt, font=\tiny] {tail $\theta_{k_i}$};
\node[note, text width=29mm] at (8.8,1.62) {new makespan, exactly:};
\node[note] at (8.8,1.18) {$M' = \max\!\big(M,\; s_j + d_j$\\[0.3ex]
  $\qquad + \max_{i \in C_j} (\tau_{j k_i} + \theta_{k_i})\big)$};
\node[note, text width=29mm] at (8.8,0.5) {scores every position without a forward pass; dashed: the arc that $p \to j \to k_i$ replaces};
% ---------------------------------------------------------------------------------------------------------------
% (c) Protocol
\draw[panel=oiGreen] (12.05,0) rectangle (17.6,5.3);
\node[head=oiGreen] at (12.1,5.25) {(c) Equal compute, pre-registered};
\node[chev=oiBlue] (s1) at (12.68,4.4) {dev:\\tune all};
\node[chev=oiOrange] (s2) at (13.98,4.4) {val:\\dry run};
\node[chev=oiGreen] (s3) at (15.28,4.4) {freeze\\(git trees)};
\node[chev=oiPink] (s4) at (16.58,4.4) {test:\\one run};
\draw[->] (s1) -- (s2); \draw[->] (s2) -- (s3); \draw[->] (s3) -- (s4);
\node[meth, draw=oiBlue, fill=oiBlue!12, line width=0.5pt] (m1) at (12.2,3.35) {\textbf{ALNS (ours)}};
\foreach \m [count=\i from 2] in {RL(s.$N$) policy, CP-SAT LNS, parallel CP-SAT LNS, CP-SAT model, CTAS-D (MILP),
                                 constructor restarts} {
  \node[meth] (m\i) at (12.2,3.72-0.37*\i) {\m};
}
\draw[decorate, decoration={brace, amplitude=3pt}] ([xshift=1mm]m1.north east) -- ([xshift=1mm]m7.south east)
  node[midway, right=3pt, font=\tiny, align=left] {same 8 pinned\\cores and\\budget $B_1$;\\CPU use\\recorded};
\node[rectangle, rounded corners=1pt, draw=black!60, fill=black!4, font=\tiny, align=center, text width=50mm,
      inner sep=1.5pt] (sim) at (14.82,0.45)
  {every plan scored by the benchmark's simulator;\\paired ratios, bootstrap CIs, Holm-corrected tests};
\draw[->] ([yshift=-0.5mm]m7.south) -- (m7.south |- sim.north);
\end{tikzpicture}
}
\caption{Overview. (a) Robots with binary trait vectors must gather in coalitions that cover each task's requirement
vector; a task starts when its last member arrives, so a late member delays the whole coalition and everything after
it. (b) Our ALNS searches over a global task order with a minimal covering coalition per task, which cannot deadlock
and is evaluated exactly; insertions are scored in time linear in the coalition size from longest-path tails
(Eq.~\eqref{eq:insert}). (c) Every method runs on the same pinned cores for the same budget, under a pre-registration
frozen before the single run on the test split.}
\label{fig:overview}
\end{figure*}

Teams of heterogeneous robots are deployed for search and rescue, inspection and construction, where many tasks need
capabilities that no single robot has \cite{kashino2020aerial,korsah2013comprehensive}. A task can then only start
once a coalition whose combined capabilities cover its requirements has gathered at its location, and every member
waits for the last one to arrive. Allocating coalitions, ordering each robot's tasks and timing the joint starts are
therefore one coupled problem: a robot that arrives late delays every task downstream of its coalition partners, and
an inconsistent ordering of shared tasks can deadlock the team. Such problems, where the schedules of different robots
constrain each other, are the hardest classes of the multi-robot task allocation taxonomies
\cite{gerkey2004formal,korsah2013comprehensive}.

Recent work has approached this problem with learned policies. \citet{dai2025heterogeneous} train a decentralised
policy by reinforcement learning (RL), with which each robot chooses its next task through an attention mechanism,
and report that it closely matches or outperforms heuristic and mixed-integer programming baselines while being at
least two orders of magnitude faster. Their best variant samples ten rollouts, RL(s.10), and scales to 150 robots and
500 tasks on the HeteroMRTA benchmark that they release, while the exact baseline, the MILP CTAS-D
\cite{fu2023robust}, solves few instances of most 50-task settings within an hour. Related learned approaches form
coalitions and routes for homogeneous teams \cite{dai2024dynamic} or schedule heterogeneous coalitions with precedence
constraints by imitation learning \cite{bichler2025sadcher}. In the HeteroMRTA evaluation, the classical competitors
are an exact solver stopped at a time limit, ant colony heuristics \cite{babincsak2023ant} and a greedy rule, and
each method runs at its own time budget.

Operations research has a long record of strong anytime heuristics for closely related problems, in particular
large neighbourhood search (LNS) for vehicle routing with synchronisation constraints
\cite{drexl2012synchronization,hojabri2018large,parragh2018solving}. It is therefore natural to ask how much of the
reported advantage of learned policies survives a comparison with a well-engineered search method and strong
constraint-programming baselines, when every method gets the same wall-clock time on the same cores. We answer this
question on HeteroMRTA with a pre-registered, compute-matched study (Figure~\ref{fig:overview}).

\paragraph{Contributions.}
\begin{itemize}
\item \textbf{An ALNS for coalition formation and scheduling with synchronisation} (Section~\ref{sec:method}). Plans
  are a global task order with one minimal covering coalition per task. This representation is deadlock-free by
  construction, its forward pass reproduces the benchmark's own simulator bit for bit, and an insertion move is
  evaluated exactly in time linear in the coalition size from longest-path tails.
\item \textbf{A compute-matched evaluation protocol} (Section~\ref{sec:setup}). All methods run in the same lanes of
  pinned physical cores at the benchmark's own budgets, with their CPU use recorded and host throttling monitored.
  Competitors are tuned on a separate split. The hypotheses, competitors, budgets and analysis were frozen, and the
  frozen code identified by its git trees, before the single run on the test split.
\item \textbf{Strong competitors.} Besides the published RL policy and constructor restarts, we compare with
  CP-SAT LNS in a threaded and a properly parallel form, with CP-SAT's own portfolio on the full model, and with
  CTAS-D, which we port to an open solver and check against the benchmark authors' own Gurobi solutions.
\item \textbf{Findings} (Section~\ref{sec:results}). ALNS is significantly better than every competitor on
  \ConeBheld{} of the 8 settings with 50 to 500 tasks, with plans \ConeBclasslo--\ConeBclasshi\% shorter than
  those of the classical baselines and \ConeBrlgainlo--\ConeBrlgainhi\% shorter than those of the learned policy.
  Parallelising CP-SAT does not close the gap, the exact MILP rarely produces a plan within the budget, and ALNS on a
  single core for 2\,s is already significantly better than every eight-core competitor at the full budget on
  \Ctwoheld{} of the 8 settings.
\end{itemize}

We take no position against learning for multi-robot coordination. Our results say that on this benchmark, at equal
compute, careful search is the stronger method, and that learned planners should be evaluated against such search
at matched budgets. The code and every result row are provided with the paper.
